% Options for packages loaded elsewhere
\PassOptionsToPackage{unicode}{hyperref}
\PassOptionsToPackage{hyphens}{url}
\documentclass[
  ignorenonframetext,
]{beamer}
\newif\ifbibliography
\usepackage{pgfpages}
\setbeamertemplate{caption}[numbered]
\setbeamertemplate{caption label separator}{: }
\setbeamercolor{caption name}{fg=normal text.fg}
\beamertemplatenavigationsymbolsempty
% remove section numbering
\setbeamertemplate{part page}{
  \centering
  \begin{beamercolorbox}[sep=16pt,center]{part title}
    \usebeamerfont{part title}\insertpart\par
  \end{beamercolorbox}
}
\setbeamertemplate{section page}{
  \centering
  \begin{beamercolorbox}[sep=12pt,center]{section title}
    \usebeamerfont{section title}\insertsection\par
  \end{beamercolorbox}
}
\setbeamertemplate{subsection page}{
  \centering
  \begin{beamercolorbox}[sep=8pt,center]{subsection title}
    \usebeamerfont{subsection title}\insertsubsection\par
  \end{beamercolorbox}
}
% Prevent slide breaks in the middle of a paragraph
\widowpenalties 1 10000
\raggedbottom
\AtBeginPart{
  \frame{\partpage}
}
\AtBeginSection{
  \ifbibliography
  \else
    \frame{\sectionpage}
  \fi
}
\AtBeginSubsection{
  \frame{\subsectionpage}
}
\usepackage{iftex}
\ifPDFTeX
  \usepackage[T1]{fontenc}
  \usepackage[utf8]{inputenc}
  \usepackage{textcomp} % provide euro and other symbols
\else % if luatex or xetex
  \usepackage{unicode-math} % this also loads fontspec
  \defaultfontfeatures{Scale=MatchLowercase}
  \defaultfontfeatures[\rmfamily]{Ligatures=TeX,Scale=1}
\fi
\usepackage{lmodern}
\ifPDFTeX\else
  % xetex/luatex font selection
\fi
% Use upquote if available, for straight quotes in verbatim environments
\IfFileExists{upquote.sty}{\usepackage{upquote}}{}
\IfFileExists{microtype.sty}{% use microtype if available
  \usepackage[]{microtype}
  \UseMicrotypeSet[protrusion]{basicmath} % disable protrusion for tt fonts
}{}
\makeatletter
\@ifundefined{KOMAClassName}{% if non-KOMA class
  \IfFileExists{parskip.sty}{%
    \usepackage{parskip}
  }{% else
    \setlength{\parindent}{0pt}
    \setlength{\parskip}{6pt plus 2pt minus 1pt}}
}{% if KOMA class
  \KOMAoptions{parskip=half}}
\makeatother
\setlength{\emergencystretch}{3em} % prevent overfull lines
\providecommand{\tightlist}{%
  \setlength{\itemsep}{0pt}\setlength{\parskip}{0pt}}
% Auto-generated by infrastructure/rendering/_pdf_latex_helpers.py::extract_math_font_preamble.
% Minimal header passed to Pandoc via -H so Beamer slides pick up the same math font
% as the combined-PDF path. Adjust the manuscript's preamble.md to change the font globally.
\usepackage{unicode-math}
\setmathfont{latinmodern-math.otf}

% Slide layout defaults for warning-clean scientific decks.
\usepackage{etoolbox}
\IfFileExists{xurl.sty}{\usepackage{xurl}}{}
\IfFileExists{seqsplit.sty}{\usepackage{seqsplit}}{\newcommand{\seqsplit}[1]{#1}}
\protected\def\breakseq#1{\seqsplit{#1}}
\protected\def\breaktt#1{\begingroup\ttfamily\seqsplit{#1}\endgroup}
\setlength{\emergencystretch}{6em}
\tolerance=5000
\hbadness=10000
\hfuzz=1pt
\setlength{\tabcolsep}{2pt}
\AtBeginEnvironment{longtable}{\tiny\renewcommand{\arraystretch}{0.86}\setlength{\tabcolsep}{1pt}}
\AtBeginEnvironment{tabular}{\tiny\renewcommand{\arraystretch}{0.86}\setlength{\tabcolsep}{1pt}}

% Natbib and cross-reference fallbacks — slides don't load natbib
% or cleveref, but combined-PDF manuscript prose may emit these
% commands. The fallback renders citations as a bracketed key list
% and cross-references as detokenized labels so slides don't fail on
% undefined control sequences or raw underscores. \providecommand is
% a no-op if packages load later.
\providecommand{\citep}[1]{[#1]}
\providecommand{\citet}[1]{#1}
\providecommand{\citealp}[1]{#1}
\providecommand{\citeauthor}[1]{#1}
\providecommand{\citeyear}[1]{#1}
\providecommand{\cref}[1]{\texttt{\detokenize{#1}}}
\providecommand{\Cref}[1]{\texttt{\detokenize{#1}}}
\usepackage{bookmark}
\IfFileExists{xurl.sty}{\usepackage{xurl}}{} % add URL line breaks if available
\urlstyle{same}
% fallback for those not using the hyperref driver hyperxmp:
\makeatletter
\@ifundefined{xmpquote}{\newcommand{\xmpquote}[1]{#1}}{}
\makeatother
\hypersetup{
  hidelinks,
  pdfcreator={LaTeX via pandoc}}

\author{\texorpdfstring{}{}}
\date{}

\begin{document}

\section{Conclusion}\label{sec:conclusion}

\begin{frame}[allowframebreaks]{Conclusion}
This study demonstrated a complete, reproducible
exploratory-data-analysis pipeline driven from a computational notebook
backed by a tested library. It validates a simple proposition: a
notebook stays trustworthy when its cells call tested code instead of
carrying logic.
\end{frame}

\begin{frame}[fragile,allowframebreaks]{Exemplar achievements}
\protect\phantomsection\label{exemplar-achievements}
Operating as the EDA exemplar for the Research Project Template
methodology, the project deployed the three foundational pillars:

\begin{enumerate}
\tightlist
\item
  \textbf{\texttt{src/eda/} library}: pure pandas/numpy data transforms
  --- loading, cleaning, summary statistics, correlation, and
  figure-data preparation --- with no plotting, no file I/O, and no
  \texttt{infrastructure} imports.
\item
  \textbf{\texttt{tests/} integrity}: a zero-mock suite over the shipped
  dataset and hand-built frames, plus a structural notebook-binding
  check, all under a ≥90\% project coverage gate.
\item
  \textbf{\texttt{docs/} knowledge operations}: architecture, testing
  philosophy, and operational rules that keep the notebook, script, and
  manuscript aligned.
\end{enumerate}
\end{frame}

\begin{frame}[fragile,allowframebreaks]{Technical contributions}
\protect\phantomsection\label{technical-contributions}
\begin{block}{The notebook -\textgreater{} tested src extraction
workflow}
\protect\phantomsection\label{the-notebook---tested-src-extraction-workflow}
The hallmark of this exemplar is the discipline it teaches: explore fast
in a cell, and the moment a computation matters, move it into
\texttt{src/eda/} behind a failing test. The walkthrough notebook
imports only the library's public surface and defines no functions of
its own, a property the test suite enforces structurally so it cannot
regress.
\end{block}

\begin{block}{Honest handling of imperfect data}
\protect\phantomsection\label{honest-handling-of-imperfect-data}
The loader surfaces missing values as \texttt{NaN} instead of silently
imputing them, and \texttt{clean\_dataset()} removes incomplete rows
with an explicit count. This makes data quality a visible, testable
property of the first pass rather than a hidden assumption.
\end{block}
\end{frame}

\begin{frame}[fragile,allowframebreaks]{Key insights}
\protect\phantomsection\label{key-insights}
\begin{enumerate}
\tightlist
\item
  \textbf{Reproducibility follows from testing fidelity}: every number
  in \breakseq{[@sec:results]} is produced by a tested function and regenerated
  on demand, so prose and artifacts cannot drift.
\item
  \textbf{Purity enables reuse}: because \texttt{src/eda/} returns
  plot-ready data and never touches a display backend, the same
  functions serve the notebook, the headless script, and the manuscript.
\item
  \textbf{Missingness is information}: reporting dropped rows beats
  imputing them away in an exploratory pass.
\end{enumerate}
\end{frame}

\begin{frame}[fragile,allowframebreaks]{Future extensions}
\protect\phantomsection\label{future-extensions}
This foundation could be extended to:

\begin{itemize}
\tightlist
\item
  \textbf{Richer cleaning}: typed imputation strategies behind tested
  functions, with the report recording which strategy ran.
\item
  \textbf{More figure preparers}: box plots, pair plots, and per-group
  overlays --- each a tested data preparer plus a thin plotting call.
\item
  \textbf{Larger / external datasets}: swap the shipped CSV for a real
  dataset by pointing \breaktt{load\_dataset(path=...)} at it while
  keeping the schema contract.
\end{itemize}
\end{frame}

\begin{frame}[fragile,allowframebreaks]{Final assessment}
\protect\phantomsection\label{final-assessment}
The \breaktt{template\_eda\_notebook} tree is the canonical reference for
how an exploratory notebook, a thin analysis script, a tested library,
and a manuscript stay synchronized across rebuilds. The pipeline
produced the figures referenced in \breakseq{[@sec:results]}, wrote
\breaktt{output/data/summary\_statistics.csv}, and rendered this markdown
together with \texttt{config.yaml} into PDF.
\end{frame}

\end{document}
